O câncer de próstata é uma das neoplasias mais comuns em homens - em 2022, foi a quarta neoplasia mais frequente no mundo, representando 7,3\% de todos os casos \cite{JournalStatistic2024}. O diagnóstico histopatológico baseia-se no sistema de graduação proposto pela \textit{International Society of Urological Pathology} (ISUP), derivado do escore de Gleason e organizado em cinco grupos (1 a 5) de progressão tumoral \cite{ref-epstein2016}. A avaliação manual de imagens de lâminas inteiras (\textit{Whole-Slide Images} - WSIs) é exaustiva e sujeita a variabilidades inter e intraobservador, o que compromete o manejo clínico \cite{ref-bulten2020}. Nesse contexto, as Redes Neurais Convolucionais (RNCs) têm representado o \textit{estado-da-arte} para automatizar a análise histopatológica \cite{ref-campanella2019}. Contudo, trabalhos anteriores frequentemente ignoram a natureza ordinal do sistema ISUP, utilizam funções de perda inadequadas para dados desbalanceados e não tratam ruídos inerentes ao processo de rotulação \cite{Xiang2023Prostate,afifi2024prostate,alici2024efficient}.

Este trabalho propõe um pipeline para a classificação automática dos graus ISUP (0-5) em WSIs prostáticas. As principais contribuições são: (i) uma função de perda híbrida que combina penalização ordinal e \textit{Focal Loss}, preservando a hierarquia biológica da doença e mitigando o desbalanceamento de classes; (ii) um método de redução de ruído baseado na entropia de Shannon das predições de um conjunto de modelos EfficientNet; e (iii) uma avaliação comparativa das arquiteturas EfficientNet-B0, B3 e B7.